\vspace{-0.15cm}
\section{Ethics Statement}
\vspace{-0.15cm}
This research involves training the \ourmethod{} model on a dataset of 17,839 publicly available arXiv pre-prints to synthesize scientific insights. While automating scientific ideation offers significant potential for discovery, it raises important ethical considerations regarding the proper attribution of foundational ideas, especially since academic citation graphs can be noisy and may not always reflect true conceptual influence. Additionally, there is an inherent risk of language models generating plausible but unverified scientific claims; our framework mitigates this by using reinforcement learning to directly align model outputs with grounded, human-derived insights. Finally, we acknowledge the environmental and computational costs associated with training models on high-performance GPUs, though utilizing a more efficient 4B parameter architecture helps limit this footprint compared to massive proprietary models.


\vspace{-0.15cm}
\section{Reproducibility Statement}
\vspace{-0.15cm}
To ensure the reproducibility of our results, we provide a comprehensive account of our methodology, code, and data. The source code for our models and experiments is available in the supplementary materials and at the following repository: \url{https://github.com/} and a subsample of our benchmark and model weights can be found in the following huggingface respository: \url{https://huggingface.co/giants2026}. Our implementation is built upon the verl framework (\url{https://verl.readthedocs.io/en/latest/}). All experimental details, including hyperparameter settings, are documented in Appendix~\ref{app:hyperparameters}. The computational experiments were conducted on a machine with NVIDIA A100 GPUs, and the required software dependencies are listed in the requirements.txt file within our code repository. The datasets used in our experiments will be made publicly available. We include samples of the dataset in the Appendix along with specific splits and any preprocessing steps applied to the data.